\subsection{Müslümanlara Uygulanan Ceza Hukuku}
\indent\indent Tanzimat öncesi Osmanlı ceza hukukunda müslümanlar şer`î hukuk ve örfî hukuk birlikte kullanılmaktaydı. Bu iki hukukun ortaklaşa işleyişini
anlamak için suçlara ve cezalara birlikte bakmak faydalı olacaktır. Şer`î hukukta suçlar genel olarak had ve tâzir olmak üzere ikiye ayrılmaktaydı. Had,
adam öldürme ve vücut bütünlüğüne yönelik işlenen suçlardı. Bunların dışında kalan suçlar ise tâzir kapsamında değerlendirilmekteydi.\footcite[99]{osmanli_hukuku}

Şer`î hukukta had suçlarının cezaları kısastır. Örfî hukuk ise had cezalarına uygulanacak para cezlarını belirlemiştir. Fatih Kannunamesi'nde bu durumla alakalı olarak şu ifade yer almaktadır:
\begin{quote}
Eğer adam öldürse yerine kısas etmeseler kan cürmi bay olup bin akçeye dahi ziyadeye gücü yeterse dört yüz, eğer altı yüze gücü yeterse iki yüz akçe, ondan aşağı halli olursa yüz akçe ve fakir olursa elli akçe alına.\footcite[101]{osmanli_hukuku}
\end{quote}

Hırsızlık suçunda ise kısas gerektiren hırsızlıklarda şer`î hukuk temel alınırken, para cezası gerektiren hırsızlıkların cezası örfî hukuk tarafından belirleniyordu. Bu noktada hırsızlık suçlarına uygulanan 
cezaların, hırsızlığın nitekli ya da basit olmasına göre değiştiği görülmektedir. Şer`î hukukta hırsızlığın nitekli olup olmaması, miktarına göre belirlenmiştir. Hanefî hukukuna göre bu miktar 10 dirhem gümüştür.
Kānunnâme-i Sultan Selim Han'da şöyle bir hükme rastlıyoruz:
\begin{quote}
Eğer bir kimesne kaz ya tavuk ya ördek uğurlasa kadı tâ'zir edüb iki ağaca bir akçe cürüm alına. Eğer bir kimesne kovan ya koyun ya kuzu uğurlasa nisabına yetişmese yine tâ'zir edüb ağaç başına bir akçe cürüm alına. Eğer bir kişi at ya katır ya merkep, uğurlasa elin keseler yahud iki yüz akçe cürüm alına.\footcite[105]{osmanli_hukuku}
\end{quote}

Bu noktada kürek cezasına değimenin de faydası olacaktır. Kürek cezası, şer`î hukukta yer almayan ancak örfî hukukun ortaya koyduğu bir ceza türüdür. Teşkilatlanmanın henüz beylik aşamasında olduğu, 14. yüzyıl 
gibi erken bir dönemde Kara Mürsel Bey'in iştiraki ile donanma faaliyetleri başlamış olsada, gerçek anlamda bir donanma teşkilatlanması 16.yüzyılda vuku bulacaktır. Büyüyen teşklatlanma ile de donanma için gerekli 
olan insan gücünün sağlanması gerekecektir. 16.yüzyıl itibariyle şer`î hukukun getirdiği kısas cezaları daha sık şekilde kürek cezasına çevrilmeye başlanacaktır.\footcite[110]{osmanli_hukuku}

Bu bilgiler ışığında, şer`î hukuk ve örfî hukukun birbiriyle çatışan değil, birbirini tamamlayan unsurlar olduğu açıktır. Örfî hukuk, şer`î hukukun getirdiği suç tanımlarına müdahil olmazken, cezaların tayiyininde 
devreye girmiştir. Örfî hukuk, had ve tâzir suç tanımlarını açıklığa kavuştururken, para cezalarının düzenlenmesinde rol oynamıştır. Değişen zaman ve ekonomik şartlar göz önünde bulundurulduğunda, bu düzenlemelerin 
siyasi otorite için bir tercihten çok bir zorunluluk olduğu anlaşılmaktadır.